%% =============================================================================
%%  config.tex — tudo o que você personaliza fica aqui
%% =============================================================================

% ---- identificação do caderno ------------------------------------------------
\newcommand\ebTitulo{Lucas 6:20–26}
\newcommand\ebSubtitulo{Sermão da Planície: felicidades e ais}
\newcommand\ebIncipit{Μακάριοι οἱ πτωχοί}
\newcommand\ebSerie{Cadernos de Estudo Bíblico}
\newcommand\ebNumero{1}
\newcommand\ebAutor{Isaac Kosloski}
\newcommand\ebData{Setembro de 2026}
\newcommand\ebEpigrafe{Those seeking Jehovah will praise him. May you enjoy life forever.}
\newcommand\ebEpigrafeRef{Psalm 22:26 — July 2024}

% ---- aparência ----------------------------------------------------------------
% Tema de cores: lapis (lápis-lazúli, padrão) | purpura | oliveira | sepia
\newcommand\ebTema{lapis}
% Papel: a4 (margem larga para notas) | b5 (formato de livro)
\newcommand\ebPapel{a4}
% Fundo levemente pergaminho em todas as páginas (bom na tela; ruim p/ imprimir)
\newcommand\ebFundoPagina{nao}
% Abas coloridas na borda da página, uma por etapa
\newcommand\ebAbas{sim}

% ---- interlinear --------------------------------------------------------------
% Modo padrão: leitura (original + translit. + pt) | estudo (+ en + morfologia)
%              | completo (+ lema e Strong)
\newcommand\ebInterlinearModo{estudo}
% Hebraico com acentos massoréticos (te'amim)? sim | nao
\newcommand\ebAcentosHebraicos{nao}

% ---- fontes alternativas (opcional) ----------------------------------------------
% Coloque os arquivos em fontes/ e descomente. Exemplos (todas livres, SIL OFL):
%   Cardo (grego e hebraico num só estilo, clássico):
%     \renewfontfamily\EBfgrego{Cardo}[Path=fontes/, Extension=.ttf,
%       UprightFont=*-Regular, ItalicFont=*-Italic, BoldFont=*-Bold]
%   Ezra SIL ou Taamey Frank CLM (hebraico tradicional, ótimos com te'amim):
%     \renewfontfamily\EBfhebraico{SILEOT}[Path=fontes/, Extension=.ttf, Script=Hebrew]
%   Crimson Pro / Libertinus Serif no lugar de EB Garamond (texto):
%     \babelfont{rm}[Path=fontes/, ...]{LibertinusSerif}
% (SBL BibLit e SBL Hebrew são gratuitas, mas com licença própria; Brill é
%  gratuita só para uso não comercial.)
\newcommand\ebFontesPersonalizadas{}
